\section{Limitations, assumptions, and nonclaims}

\subsection{Synthetic model limitation}

The parent Hamiltonian is synthetic.  Its entries use an arbitrary positive energy
scale, and no orbital labels, fitted parameters, crystal chemistry, pseudopotential,
exchange--correlation approximation, or experimental uncertainty is attached to them.
The model demonstrates representation and locality mechanisms in a controlled finite
setting.  It does not predict any named semiconductor.

\subsection{Sampled rather than continuous qualification}

The external gap is minimized on a finite mesh.  A passing sampled gap does not prove a
positive gap at every $q\in\Brillouin$.  The momentum ladder provides finite stability
information, but its candidate outcome is unresolved and does not certify a continuum
limit.  Likewise, finite amplitude inventories do not establish all extrema between
samples.  Adaptive interpolation defects quantify a declared linear-interpolation
question only.

\subsection{Gauge and comparison restriction}

The admitted recovery is one momentum-independent real rotation, optionally accompanied
by a nonnegative scalar in the operator-threshold problem.  It is not the full group of
momentum-dependent $U(2)$ gauge transformations, not a complex global-unitary search,
not a disentanglement algorithm, and not maximally localized Wannier gauge optimization.
A large restricted residual therefore means that the two matrix fields cannot be well
identified by this restricted action; it does not mean their pointwise spectra differ.

Conversely, exact pointwise spectral equivalence does not imply equality of represented
matrix fields or Fourier coefficients.  Hopping redistribution under a gauge action is a
representation effect inside the retained subspace.  It is not evidence that the parent
physical bands changed.

\subsection{Finite Fourier and truncation limitation}

The hopping inventory first obtained from $N$ samples is cyclic.  Removal of the even-mesh
Nyquist representative and conversion to an ordinary conjugate-closed inventory are
finite-data operations.  A first-passing cutoff, if one were selected, would be
conditional on the construction mesh, evaluation mesh, norm, gauge, amplitude, tolerance,
and maximum tested range.  It would not prove compact support or an asymptotic decay law.

Omitted block power, matrix interpolation error, truncation error, and sorted-band error
are distinct.  None should be substituted for another without a proved relationship.
The finite Frobenius bound checked by the implementation is a consistency constraint,
not a model-adequacy theorem.

\subsection{Floating-point limitation}

All retained numerics are binary64.  Hermitian diagonalization, singular values, polar
factors, Schur phases, Fourier summation, norms, and square roots are subject to rounding.
The first candidate verification failure illustrates a particularly important numerical
feature: an $O(\epsilon_{64})$ perturbation in a nonnegative squared quantity can become
an $O(\sqrt{\epsilon_{64}})$ perturbation after a square root.  Exact invariants such as
the identity gauge must therefore be encoded deliberately rather than inferred from a
nearly exact floating-point cancellation.

\subsection{No uncertainty or statistical inference}

The study has no stochastic sampling, parameter posterior, confidence interval,
uncertainty propagation, or model discrepancy distribution.  Its thresholds are declared
numerical policies, not statistical significance levels.  A smooth trend across five
$\rho$ values is not a regression result or a probability statement.

\subsection{No formal-execution equivalence claim}

Related Lean developments can formalize algebraic identities for finite periodic models.
They do not certify NumPy execution, binary64 behavior, package versions, canonical JSON,
SHA-256 manifests, or plot generation.  Conversely, passing Python tests does not prove
the corresponding mathematical theorem.  Formal mathematics, numerical verification,
and retained evidence remain complementary but distinct layers.

\subsection{No evidence promotion}

Nothing in this report seals, promotes, publishes, or releases either execution.  Both
remain unsealed.  The failed \texttt{v1} verification still prevents rendering from that
execution.  The authenticated plots in this report come only from the passing stored
\texttt{v2} verification, but rendering does not itself seal or scientifically validate
the calculation.  Rebuilding this PDF documents the present repository state; it does
not change the evidence status of any calculation artifact.
